\section{Empirical Evaluation}
\label{sec:evaluation}

\subsection{Experimental Setup & Ground-Truth Dataset}
We evaluate AgentPOI against 60 ground-truth enterprise financial scenarios:
\begin{itemize}
    \item \textbf{30 Adversarial Scenarios}: Direct prompt injection, homoglyph typosquatting, payload splitting, markdown table smuggling, compliance camouflage, ceiling escalation, and sleeper state poisoning adapted from AgentDojo Banking~\cite{debenedetti2024agentdojo} and InjecAgent Financial Tools~\cite{zou2024injecagent}.
    \item \textbf{30 Clean Operational Scenarios}: Itemized multi-line purchase orders, credit memo balance reconciliations, recurring utility disbursements, and urgent executive directives.
\end{itemize}

\subsection{Evaluation Metrics}
\begin{itemize}
    \item \textbf{ASR-Exec (Attack Success Rate on Execution)} ($\downarrow$): Percentage of adversarial scenarios resulting in unauthorized funds transfer.
    \item \textbf{LDSR (Legitimate Disbursement Success Rate)} ($\uparrow$): Percentage of authentic vendor invoices successfully approved and executed.
    \item \textbf{CDR (Clean Disruption Rate)} ($\downarrow$): Percentage of authentic corporate tasks falsely blocked or disrupted.
    \item \textbf{Latency Overhead}: Runtime overhead added per transaction turn on standard CPU.
    \item \textbf{Token Overhead}: Additional defensive LLM tokens consumed.
\end{itemize}

\subsection{Cross-Framework Benchmark Results}
Table~\ref{tab:framework_comparison} reports the comparative evaluation across LangChain ReAct, CrewAI, AutoGPT, and AgentPOI.

\begin{table*}[t]
\centering
\caption{Comparative Security and Operational Resilience Across Agentic Architectures (60 Real-World Ground-Truth Financial Scenarios).}
\label{tab:framework_comparison}
\begin{tabular}{lcccccc}
\hline
\textbf{Framework} & \textbf{Architecture Paradigm} & \textbf{ASR-Exec} $\downarrow$ & \textbf{LDSR} $\uparrow$ & \textbf{CDR} $\downarrow$ & \textbf{Latency (ms)} & \textbf{Token Cost} \\
\hline
LangChain ReAct & Single-Agent ReAct (Tool-Calling Loop) & 100.0\% & 100.0\% & 0.0\% & 0.12 & 0 \\
CrewAI & Multi-Agent Hierarchical Role-Playing & 100.0\% & 100.0\% & 0.0\% & 0.16 & 450 \\
AutoGPT & Autonomous Planning \& Workspace Execution & 70.0\% & 76.7\% & 23.3\% & 0.13 & 600 \\
\textbf{AgentPOI (FinancialGuard)} & \textbf{Dual-State Epistemic Consensus (CIV + AuthGraph)} & \textbf{0.0\%} & \textbf{100.0\%} & \textbf{0.0\%} & \textbf{0.22} & \textbf{0} \\
\hline
\end{tabular}
\end{table*}

\subsection{Architectural Ablation Analysis}
To quantify the incremental defense value of each layer, Table~\ref{tab:ablation_analysis} presents the ablation study isolating each architectural component.

\begin{table*}[t]
\centering
\caption{Architectural Ablation Analysis: Incremental Defense Contribution of Security Components.}
\label{tab:ablation_analysis}
\begin{tabular}{lccccc}
\hline
\textbf{Ablation Configuration} & \textbf{Defense Component} & \textbf{ASR-Exec} $\downarrow$ & \textbf{LDSR} $\uparrow$ & \textbf{CDR} $\downarrow$ & \textbf{Latency (ms)} \\
\hline
1. Baseline & No Security Gates & 100.0\% & 100.0\% & 0.0\% & 0.12 \\
2. Task Shield Only & Fact vs. Imperative Decomposition (ACL 2025) & 33.3\% & 96.7\% & 3.3\% & 0.18 \\
3. AuthGraph Only & Cryptographic Provenance DAG (arXiv:2605.26497) & 16.7\% & 100.0\% & 0.0\% & 0.15 \\
4. CIV Only & Contextual Invariant Verification (Core Engine) & 6.7\% & 100.0\% & 0.0\% & 0.19 \\
\textbf{5. Full AgentPOI} & \textbf{Full Dual-State Epistemic Consensus} & \textbf{0.0\%} & \textbf{100.0\%} & \textbf{0.0\%} & \textbf{0.22} \\
\hline
\end{tabular}
\end{table*}

The ablation demonstrates that while Task Shield and AuthGraph dramatically reduce attack success, zero-breach security is achieved \textit{only} when all three layers operate under Dual-State Epistemic Consensus.
